\section{Introduction}
\label{sec:intro}

A precision farm is a multi-stakeholder system before it is a technical one. Soil probes, weather stations and irrigation actuators produce readings continuously \citep{zhang2002precision}, and the people who consume those readings want different and partly conflicting things from them: farmers need operational control, agronomists need history but no actuation rights, certification bodies need compliance evidence without visibility into commercially sensitive agronomic detail, supply-chain partners need provenance but not real-time operations, and regulators need environmental records. Handing this authorization matrix to a single server concentrates the risk in one place. An attacker who takes that server gains not only the data but, on a farm with networked valves, the irrigation itself \citep{benedict2021decentralized}, and the audit log that would reconstruct the incident is written by the machine that was compromised. Surveying agricultural IoT deployments, Ferrag et al.\ \citep{ferrag2020security} found most lacked fine-grained access control at all, and Kamilaris et al.\ \citep{kamilaris2019rise} and recent convergence surveys \citep{rustemi2025synergizing} find blockchain work in agriculture concentrated almost entirely on post-harvest supply-chain tracking rather than pre-harvest authorization. Recording role assignments, revocations and access decisions on a permissioned ledger can reduce dependence on a single policy database and make authorized changes auditable across participating organizations. It does not eliminate trust in the certificate authority, administrators, gateway software, endorsement configuration or chaincode implementation. What is unsettled is the cost of the part it does change, and specifically where that cost falls.

The constraints behind the design in Section~\ref{sec:design} are not incidental to the measurement; they shape what the field deployment observes beyond the controlled performance campaigns. Constrained sensor energy budgets push caching and enrollment logic to the edge gateway, and intermittent rural connectivity, not injected fault, is why the revocation and outage behaviour of Sections~\ref{sec:revocation-results} and~\ref{sec:availability} includes a certificate revoked while its gateway is disconnected, a failure mode not represented in the deployed benchmark measurements. LoRa's low bit rate and eFuse-only key storage set the engineering context for the residue-coded transmission and signature paths audited in Sections~\ref{sec:crt-audit} and~\ref{sec:sig-audit}, but their defects were exposed by cross-component audit rather than deployment duration. None of this is unique to agriculture, but the combination -- long-running, energy-constrained, intermittently connected and physically exposed -- makes the field record a different operational instrument from a data-centre benchmark of the same chaincode.

Among the systems surveyed in the supplementary material, evaluation is typically a laboratory or simulated run of hours to days, reporting peak throughput and mean latency that aggregate over a pipeline of at least four separable stages, so an engineer told a transaction takes \TotalLatencyPeerFour{} ms cannot tell whether to rewrite the policy engine, restructure endorsement, replace the ordering service, or accept the figure as a consensus floor. Existing evaluations rarely combine long-duration field observations with an implementation-level authorization timing span, controlled configuration experiments and a post-deployment audit; that combination, not any single measurement, is the gap this paper addresses. Attributing time to Fabric's internal stages needs one transaction timestamped at each hand-off, which our instrumentation does not do, so we report a coarser accounting instead: the chaincode logs a timing field on every write, and because its own start and stop boundaries were not documented at deployment time we treat it as an authorization-associated span rather than a stage measurement, with endorsing-peer count acting as a lever on what remains. Abang et al.\ \citep{abang2024latency}, reviewing Fabric latency modelling for IoT, make a closely related point: the literature models latency in the abstract but rarely connects those models to a specific deployed use case measured long enough for its slow-moving parts to misbehave. We do not supply that stage-level measurement either; what we add is a long-duration field record against which the recorded span holds up, plus a code-to-specification audit that identified three cross-layer defect classes -- authorization-policy deviations, invalid residue recovery and signature-verification failure -- and one revocation-observability gap; the deployed benchmark measurements and component-level tests detected none of the three defect classes.

The intuitive expectation turns out to be wrong, which is the main reason we think the measurement is worth reporting. Access control is the component under study, so it is natural to assume it dominates the cost. The recorded span does not: it is a modest and, across the tested peer configurations, comparatively stable component of the write path, and the larger, more variable share sits in the layer that access-control papers usually treat as background infrastructure. A second gap is temporal: certificate lifecycle management, revocation propagation and enrollment recovery are the mechanisms most likely to fail in the field and least likely to be exercised by short-duration benchmark measurements, and a revocation pipeline that behaves perfectly across fifty laboratory trials can still show a materially different delay distribution across two months of intermittent rural connectivity.

These questions structure the rest of the paper:

\textbf{RQ1} What latency and throughput differences are observed between the tested access-control-enabled configuration and the access-control-free comparison configuration?

\textbf{RQ2} How is configured logical peer count associated with client-observed write latency under the deployed topology?

\textbf{RQ3} Which implementation and observability failures survived the field deployment, and why did the existing tests not detect them?

We built an HRBAC framework on Hyperledger Fabric 2.5 for the constraints of agricultural edge deployment, ran it on a working farm for \DeploymentDays{} days, and instrumented the authorization-associated span and the end-to-end transaction path. Four bounded contributions follow, each naming which evidence it rests on -- the field deployment, a later controlled laboratory campaign, or the code audit -- and none extending its claim beyond that source: the field observations characterize the deployed system as run; they are not measurements of a corrected implementation (Section~\ref{sec:corrected-absent}) or evidence of universal blockchain scalability.

\textbf{A \DeploymentDays{}-day field characterisation}, \AuthDecisions{} authorization decisions over \Sensors{} sensors, \Gateways{} gateways and \Zones{} zones, addressing RQ1 (Section~\ref{sec:results}).

\textbf{A controlled comparison across tested configurations}: the recorded authorization-associated span averages \SpanLatencyPeerFour{} ms, \SpanSharePeerFour\%\ of the \TotalLatencyPeerFour{} ms write path at 4 peers, and a counterbalanced comparison of the 4- and 32-peer configurations finds it changes by only \SpanDiffFourVThirtyTwo{} ms while the remaining latency falls by \RemainingLatencyDiffFourVThirtyTwo{} ms, addressing RQ2; neither the span's internal boundaries nor the remainder's internal stages are separately identified (Sections~\ref{sec:results-latency-accounting} and~\ref{sec:results-peer-scaling}).

\textbf{An implementation postmortem}: an audit of the implementation against its own specification found three cross-layer defect classes -- authorization-policy deviations, invalid residue recovery and signature-verification failure -- that survived \DeploymentDays{} days of operation, plus one revocation-observability gap, addressing RQ3. The deployed benchmark measurements and component-level tests did not detect the defect classes; together, the audit findings are the part of this work most likely to transfer to other deployments (Section~\ref{sec:audit}).

\textbf{A reproducible artifact and cross-layer testing recommendations}: the raw traces, chaincode, firmware and the analysis pipeline that recomputes every number, table and figure in this article are released at the repository named in the Data Availability statement (Section~\ref{sec:data-availability}). Unlike an earlier internal draft of this analysis, the corrected implementation described in the postmortem is not, at the time of this submission, present in that repository as a separate tagged package; Section~\ref{sec:corrected-absent} states this limitation directly rather than assuming the correction exists because it is specified.

Section~\ref{sec:related} positions the work. Section~\ref{sec:design} presents the access-control model and the enrollment protocol, Section~\ref{sec:implementation} the implementation, Section~\ref{sec:methodology} the measurement methodology, Section~\ref{sec:results} the results, Section~\ref{sec:discussion} what follows from them for system design, and Section~\ref{sec:conclusion} concludes.
